\begin{abstract}
Test-suite maintenance tools lean on cheap signals: tests that cover the same lines look redundant, a function whose deletion breaks a test looks protected, and a test that passed a few reruns looks stable. Before building these signals into the Test Suite Auditor, a Claude Code plugin, we tested six pre-registered hypotheses about them on eight mature Python projects (\NTestsInScope{} in-scope tests), using the plugin's own runner. Among test pairs with coverage Jaccard of at least 0.9, \HOneShare{} (95\% CI \HOneLo{}--\HOneHi{}) are complementary: each kills a mutant the other misses. Of the call deletions that no test noticed, \HThreeShare{} (\HThreeLo{}--\HThreeHi{}) sit inside functions whose every whole-body mutant was killed, a stronger version of a result previously reported for Java. Removing all individually redundant tests from a file together lost kills in \HFourLosing{} of \HFourEligible{} files. Per-test coverage contexts slowed suites by a pooled factor of \HFiveRatio{}, and the only intermittent tests we observed in 60 runs per project were order-dependent: in file order one fails every time and the other never fails. \HTwoAbstract{} We release the instrument, raw data and analysis, and describe how each result changes the plugin's design.
\end{abstract}
